\section{The Testbed and Why Its Numbers Are Trustworthy}
\label{sec:testbed}

The model is fully discrete, so it is fully simulable. The testbed is an open-source
TypeScript library (the \texttt{@cluesurf/vibe} package). It generates a substrate,
runs a local rule over it in discrete beats, and measures the emergent physics.
Everything is finite and seeded, so every number in this paper is a deterministic
function of a stated seed and reproduces exactly. Real numbers appear only as measured
outputs (coordinates, eigenvalues), never as the base.

At the core is a suite of \emph{known-answer tests}. A number on its own means
nothing. A number from a tool checked against a case with an independently known answer
does. Every measurement tool is calibrated against such a case: a Poisson sprinkling of
a $d$-dimensional region must measure dimension $d$, a Bell test with independent
settings must respect the classical bound of $2$, the overlap operator must have exactly
one fermion species with machine-zero chiral-symmetry residual, the uniform-measure
sampler must reproduce exact enumeration, and so on. There are fifty-six such checks,
and they all pass. Several early versions were wrong in plausible-looking ways, a
miscalibrated dimension estimator, a sprinkler not uniform by proper volume, a Monte
Carlo move that did not sample the right measure, and the known-answer tests caught them.
A program that cannot be caught being wrong is not science. This one can be, and was,
repeatedly.

What the testbed establishes falls in three honest categories. First, the machinery is
correct, validated by the known-answer tests. Second, the framework can carry known
physics: reproducing the Newtonian potential, the photon, the graviton, confinement,
and the electroweak masses shows the mesh is expressive enough to hold the real world,
which is a load-bearing test, since a monist theory is falsified if matter, force, and
gravity cannot be built from one substrate. It is not, by itself, proof the framework
is uniquely correct, since standard physics reproduces these too. Third, the genuinely
distinctive claims, where the framework makes a stand of its own: the both-worlds
substrate, the dynamical selection of spacetime, the superdeterministic quantum
tension, the Lorentz-safety of random discreteness, and the everpresent cosmological
constant. The distinctive predictions that meet data (Section~\ref{sec:predictions})
are what could confirm or falsify the theory against nature.
